% transformer_slides.tex -- additional Transformer-style surrogate slides
% Source context: MIKU-7437/AI_Termoelastic_Waves_Geology, chapter 5 Transformer file.

\begin{frame}{Transformer: Attention-Based Surrogate}
  \begin{columns}[T]
    \begin{column}{0.50\textwidth}
      \begin{block}{Role in the framework}
        \footnotesize
        The Transformer-style route is an attention-based surrogate for fast
        prediction under the shared input--output contract.
      \end{block}
      \begin{itemize}
        \footnotesize
        \item Treats sampled points as tokens.
        \item Relates distant points through attention.
        \item Supports query-style prediction at requested locations.
      \end{itemize}
    \end{column}
    \begin{column}{0.46\textwidth}
      \begin{alertblock}{Interpretation}
        \footnotesize
        This route is mainly data-driven. It can be fast and flexible, but its
        physical priors are weaker than those of the PINN route.
      \end{alertblock}
    \end{column}
  \end{columns}
\end{frame}

\begin{frame}{Transformer Tokenisation}
  \begin{columns}[T]
    \begin{column}{0.52\textwidth}
      \begin{block}{Input token}
        \[
          a_i=[x_i,y_i,T_i^{(0)},u_i^{(0)},v_i^{(0)},
          \dot u_i^{(0)},\dot v_i^{(0)},E_i,\nu_i,\rho_i,
          \alpha_i,k_i,C_{p,i}]
        \]
      \end{block}
    \end{column}
    \begin{column}{0.44\textwidth}
      \begin{block}{Query token}
        \[
          q_j=(x_q^{(j)},y_q^{(j)})
        \]
      \end{block}
      \footnotesize
      Input tokens describe known mesh samples. Query tokens describe where the
      predicted fields are requested.
    \end{column}
  \end{columns}
\end{frame}

\begin{frame}{Scaled Dot-Product Attention}
  \begin{block}{Attention operation}
    \[
      \operatorname{Attention}(Q,K,V)=
      \operatorname{softmax}\left(\frac{QK^{\top}}{\sqrt{d_k}}\right)V
    \]
  \end{block}
  \begin{columns}[T]
    \begin{column}{0.48\textwidth}
      \footnotesize
      \begin{itemize}
        \item $Q$ stores query representations.
        \item $K$ stores key representations.
        \item $V$ stores value representations.
      \end{itemize}
    \end{column}
    \begin{column}{0.48\textwidth}
      \footnotesize
      The score matrix connects every query with every input token. This gives
      the model a direct route for long-range interactions in the domain.
    \end{column}
  \end{columns}
\end{frame}

\begin{frame}{Transformer Encoder and Decoder}
  \begin{columns}[T]
    \begin{column}{0.48\textwidth}
      \begin{block}{Encoder}
        \footnotesize
        Input tokens are projected to an embedding space and refined through
        multi-head self-attention and feed-forward layers.
      \end{block}
      \[
        H^{(\ell+1)}=\operatorname{EncoderBlock}(H^{(\ell)})
      \]
    \end{column}
    \begin{column}{0.48\textwidth}
      \begin{block}{Decoder}
        \footnotesize
        Query tokens attend to encoded input tokens through cross-attention.
      \end{block}
      \[
        Z^{(\ell+1)}=\operatorname{DecoderBlock}(Z^{(\ell)},H)
      \]
    \end{column}
  \end{columns}
\end{frame}

\begin{frame}{Transformer Output Head}
  \begin{columns}[T]
    \begin{column}{0.52\textwidth}
      \begin{block}{Predicted fields at query points}
        \[
          \hat{y}_j=(\hat T_j,\hat u_j,\hat v_j)
        \]
        \[
          \hat{y}_j=W_{\mathrm{out}}Z_{j,:}^{(L_{\mathrm{dec}})}
        \]
      \end{block}
    \end{column}
    \begin{column}{0.44\textwidth}
      \footnotesize
      The model predicts temperature and two in-plane displacement components.
      Stress and strain can be derived later from displacement gradients or
      added as future output channels.
      \vskip 4pt
      \scriptsize
      The baseline described in the source text uses an embedding dimension of
      $d_m=128$ and multi-head attention with $h=4$ heads.
    \end{column}
  \end{columns}
\end{frame}

\begin{frame}{Transformer Training Objective}
  \begin{block}{Hybrid supervised objective}
    \[
      \mathcal{L}=\lambda_{\mathrm{sup}}\mathcal{L}_{\mathrm{sup}}
      +\lambda_{\mathrm{vel}}\mathcal{L}_{\mathrm{vel}}
      +\lambda_{\mathrm{th}}\mathcal{L}_{\mathrm{th}}
    \]
  \end{block}
  \begin{columns}[T]
    \begin{column}{0.48\textwidth}
      \footnotesize
      \begin{itemize}
        \item supervised reconstruction of reference fields;
        \item velocity consistency for displacement evolution;
        \item thermal residual regularisation.
      \end{itemize}
    \end{column}
    \begin{column}{0.48\textwidth}
      \footnotesize
      Token subsampling controls attention cost and encourages operation on
      different spatial samplings. Predictions still require careful validation
      against COMSOL-derived references.
    \end{column}
  \end{columns}
\end{frame}
